\documentclass[10pt,twocolumn]{article}

% ── packages ──────────────────────────────────────────────────────────────────
\usepackage[utf8]{inputenc}
\usepackage[T1]{fontenc}
\usepackage{amsmath,amssymb}
\usepackage{graphicx}
\usepackage{booktabs}
\usepackage{hyperref}
\usepackage[margin=0.75in]{geometry}
\usepackage{natbib}
\usepackage{xcolor}
\usepackage{algorithm}
\usepackage{algpseudocode}
\usepackage{enumitem}
\usepackage{microtype}
\tolerance=1500
\emergencystretch=1em
\usepackage{multirow}
\usepackage{subcaption}
\usepackage{tikz}
\usetikzlibrary{arrows.meta,positioning,calc}

\hypersetup{colorlinks=true, linkcolor=blue!60!black, citecolor=blue!60!black, urlcolor=blue!60!black}

\title{Loom v2: From Shared Embeddings to Task-Specific Heads\\for Outfit Recommendation with\\Object-Centric Wardrobe Understanding}

\author{
  Anushree Berlia
}

\date{May 2026}

\begin{document}
\maketitle

% ══════════════════════════════════════════════════════════════════════════════
\begin{abstract}
We present \textbf{Loom~v2}, a substantial evolution of the Loom outfit
recommendation system~\citep{berlia2026loom} that replaces a single shared
FashionCLIP embedding with a \emph{multi-head projection architecture}
built on a frozen DINOv2 backbone.  Five task-specific MLP heads produce
128-dimensional embeddings for style, fit, material, compatibility, and
occasion---each trained via contrastive learning on Polyvore, DeepFashion,
and iMaterialist data.  A learned compatibility scorer replaces the
rule-based four-level outfit hierarchy with pairwise neural scoring.

Beyond the representation layer, Loom~v2 introduces three additional
architectural advances:
(1)~\emph{object-centric wardrobe ingestion} via YOLOv8-Nano detection and
ByteTrack temporal tracking, enabling multi-item capture from video and
live camera feeds;
(2)~\emph{cascade retrieval} with three-round sequential conditioning
(silhouette anchoring, third-piece context, finisher completion); and
(3)~\emph{segmentation-first embedding} using U\textsuperscript{2}-Net
background removal before all encoders.

Ablation experiments on a 620-item catalog show that the multi-head system
achieves a mean outfit score of 0.489 with 0\% hard violations, compared
to 0.011 score and 8.0\% violations for the rule-based baseline---a
44$\times$ improvement in score and complete elimination of constraint
violations.  The system generates three stylistically distinct outfits in
under 5.1 seconds on commodity hardware, 22\% faster than the
FashionCLIP-only predecessor.

\smallskip\noindent
Code: \url{https://github.com/anushreeberlia/loom}
\end{abstract}

% ══════════════════════════════════════════════════════════════════════════════
\section{Introduction}
\label{sec:intro}

The original Loom system~\citep{berlia2026loom} demonstrated that
combining neural embedding retrieval with structured domain scoring could
produce coherent outfits from fashion catalogs.  However, its reliance on
a single shared FashionCLIP~\citep{chia2022fashionclip} embedding for
\emph{all} downstream tasks---retrieval, occasion filtering, style
matching, material compatibility---created a fundamental tension:
FashionCLIP was trained to align images with product descriptions, not to
predict outfit compatibility or occasion appropriateness.

This mismatch manifested in several practical failures:
\begin{itemize}[leftmargin=*,topsep=2pt,itemsep=1pt]
\item Items that ``look alike'' (high cosine similarity) were retrieved for
  slots where \emph{complementarity} was needed, not similarity.
\item Occasion filtering required hand-crafted ``vibe/anti-vibe'' prose
  strings with manually tuned penalty coefficients
  ($\lambda_{\text{work}} = 3.0$).
\item Material compatibility checking via CLIP heavy/light context vectors
  was fragile, rejecting valid combinations and narrowing candidate pools.
\item The four-level hierarchical scorer required hundreds of lines of
  rule-based code for silhouette, color, texture, and finishing
  assessments.
\end{itemize}

Loom~v2 addresses these limitations by replacing the shared embedding with
task-specific projection heads, each trained to capture a different
dimension of fashion understanding.  We also introduce three
infrastructure advances---object-centric ingestion, cascade retrieval, and
segmentation-first processing---that collectively transform Loom from a
static catalog tool into a persistent wardrobe understanding system.

\paragraph{Contributions.}
\begin{enumerate}[leftmargin=*,topsep=2pt,itemsep=1pt]
\item \textbf{Multi-head projection architecture}: Frozen DINOv2
  backbone with five contrastively trained 128-d heads and a learned
  pairwise compatibility scorer (Section~\ref{sec:multihead}).
\item \textbf{Object-centric wardrobe ingestion}: YOLOv8-Nano detection
  $+$ ByteTrack tracking for multi-item video and live camera capture
  (Section~\ref{sec:ingestion}).
\item \textbf{Cascade retrieval with sequential conditioning}:
  Three-round retrieval that builds outfits progressively---silhouette
  anchor, third piece, finisher---with each round conditioned on prior
  selections (Section~\ref{sec:cascade}).
\item \textbf{Segmentation-first embedding}: U\textsuperscript{2}-Net
  background removal before all encoders, eliminating room/body/furniture
  noise from fashion representations (Section~\ref{sec:segmentation}).
\item \textbf{Comprehensive ablation}: Head-by-head and component-level
  evaluation against the rule-based predecessor on the same catalog and
  anchor set (Section~\ref{sec:experiments}).
\end{enumerate}

% ══════════════════════════════════════════════════════════════════════════════
\section{Related Work}
\label{sec:related}

\paragraph{Self-supervised vision transformers.}
DINOv2~\citep{oquab2024dinov2} produces patch-level and CLS-level features
via self-supervised distillation.  Its representations transfer strongly to
downstream tasks without fine-tuning, making it suitable as a frozen
backbone for multi-head projection.

\paragraph{Multi-task fashion embeddings.}
Type-aware embeddings~\citep{vasileva2018learning} project items into
type-conditioned spaces.  MVAE~\citep{han2019mvae} learns disentangled
latent codes.  Our approach differs by using \emph{independent} projection
heads rather than a shared latent space, allowing each head to specialize
without cross-task interference.

\paragraph{Contrastive compatibility learning.}
SCE-Net~\citep{tan2019scenet} learns compatibility via conditional
similarity.  We use a similar contrastive formulation (InfoNCE) but
decompose compatibility into multiple orthogonal heads, each with
domain-specific training pairs.

\paragraph{Object detection in fashion.}
ModaNet~\citep{zheng2018modanet} provides fashion-specific detection.
DeepFashion2~\citep{ge2019deepfashion2} extends to landmark detection.
We use general-purpose YOLOv8-Nano~\citep{jocher2023ultralytics} for
detection and add ByteTrack~\citep{zhang2022bytetrack} temporal tracking
for video-mode multi-item capture.

\paragraph{Cascade and sequential recommendation.}
Sequential outfit composition has been explored via
autoregressive generation~\citep{chen2019pog} and iterative
refinement~\citep{lin2020csanet}.  Our cascade retrieval differs by using
\emph{fixed-stage} sequential conditioning (silhouette $\to$ third piece
$\to$ finisher) aligned with how fashion practitioners build outfits.

% ══════════════════════════════════════════════════════════════════════════════
\section{System Architecture}
\label{sec:architecture}

\begin{figure*}[t]
\centering
\fbox{\parbox{0.95\textwidth}{\centering
\textbf{Loom v2 Pipeline}\\[6pt]
\texttt{Camera/Video/Photo}
$\;\xrightarrow{\text{YOLO+ByteTrack}}\;$
\texttt{Per-Object Best Crop}
$\;\xrightarrow{\text{U\textsuperscript{2}-Net}}\;$
\texttt{Segmented Cutout}\\[4pt]
$\downarrow$\\[4pt]
\texttt{DINOv2 ViT-B/14} $\to$ \texttt{768-d backbone}
$\;\xrightarrow{\text{5 MLP heads}}\;$
\texttt{5 $\times$ 128-d task vectors}\\[4pt]
$\downarrow$\\[4pt]
$\forall\, s \in \{\text{bottom, shoes, layer, accessory}\}:$\quad
\texttt{compat\_emb query}
$\;\xrightarrow{\text{pgvector HNSW}}\;$
\texttt{Candidates}$_s$
$\;\xrightarrow{\text{occasion\_emb filter}}\;$
\texttt{Filtered}$_s$\\[4pt]
$\downarrow$\\[4pt]
\texttt{Cascade: silhouette $\to$ third piece $\to$ finisher}
$\;\xrightarrow{\text{CompatScorer}}\;$
\texttt{Best outfit per direction}
}}
\caption{Loom~v2 end-to-end pipeline.  Compared to v1, ingestion adds
object detection and segmentation; embedding replaces FashionCLIP with
DINOv2+heads; retrieval uses compatibility embeddings; scoring uses a
learned pairwise scorer.}
\label{fig:pipeline_v2}
\end{figure*}

% ══════════════════════════════════════════════════════════════════════════════
\section{Method}
\label{sec:method}

\subsection{Multi-Head Projection Architecture}
\label{sec:multihead}

\paragraph{Backbone.}
We use DINOv2 ViT-B/14~\citep{oquab2024dinov2} as a frozen visual
backbone, extracting the 768-dimensional CLS token from
$224 \times 224$ center-cropped images.  DINOv2 was chosen over FashionCLIP
because its self-supervised training objective produces richer
\emph{visual} features without conflating them with text alignment.

\paragraph{Projection heads.}
Five independent MLP heads project the backbone embedding into
task-specific 128-dimensional spaces:

\begin{equation}
  \mathbf{h}_t = \text{L2norm}\Big(
    W_2^{(t)} \cdot \text{LN}\big(
      \text{ReLU}(W_1^{(t)} \mathbf{x} + b_1^{(t)})
    \big) + b_2^{(t)}
  \Big)
\label{eq:head}
\end{equation}

\noindent where $\mathbf{x} \in \mathbb{R}^{768}$ is the DINOv2 CLS
token, $W_1^{(t)} \in \mathbb{R}^{256 \times 768}$,
$W_2^{(t)} \in \mathbb{R}^{128 \times 256}$, and $\text{LN}$ is
layer normalization.  Each head $t$ specializes in one dimension:

\begin{table}[t]
\centering
\small
\caption{Multi-head architecture: five projection heads with training
data sources and target semantics.}
\label{tab:heads}
\begin{tabular}{@{}llp{3.5cm}@{}}
\toprule
Head & Training data & Target semantics \\
\midrule
\texttt{style}    & Polyvore & Aesthetic identity (minimalist, streetwear, romantic) \\
\texttt{fit}      & DeepFashion & Silhouette/cut (oversized, tailored, bodycon) \\
\texttt{material} & iMaterialist & Fabric/texture (wool, denim, silk) \\
\texttt{compat}   & Polyvore outfits & Co-outfit compatibility \\
\texttt{occasion} & DeepFashion & Context (work, casual, going-out) \\
\bottomrule
\end{tabular}
\end{table}

\paragraph{Contrastive training.}
Each head is trained independently using InfoNCE~\citep{oord2018infonce}
with temperature $\tau = 0.07$:
\begin{equation}
  \mathcal{L}_{\text{head}} = -\log \frac{
    \exp(\mathbf{h}_i \cdot \mathbf{h}_j^+ / \tau)
  }{
    \sum_{k=1}^{N} \exp(\mathbf{h}_i \cdot \mathbf{h}_k / \tau)
  }
\label{eq:infonce}
\end{equation}

\noindent Positive pairs $(i, j^+)$ are constructed from dataset-specific
annotations: co-outfit items for compatibility, same-style items for
style, same-attribute items for material/fit/occasion.  The backbone is
frozen; only the head parameters ($\sim$200K per head, 1M total) are
trained.  Training uses AdamW (lr $10^{-3}$, weight decay $10^{-2}$) for
30 epochs (40 for compatibility with hard negative mining after epoch 10),
batch size 256, on pre-computed HDF5 backbone embeddings.

\paragraph{Compatibility scorer.}
A separate 3-layer MLP predicts pairwise compatibility from concatenated
\texttt{compat} head embeddings:
\begin{equation}
  p(a, b) = \sigma\Big(
    W_3 \cdot \text{ReLU}\big(
      W_2 \cdot \text{ReLU}(W_1 [\mathbf{h}_a^c ; \mathbf{h}_b^c] + b_1) + b_2
    \big) + b_3
  \Big)
\label{eq:compat}
\end{equation}

\noindent where $[\cdot;\cdot]$ denotes concatenation,
$W_1 \in \mathbb{R}^{128 \times 256}$, $W_2 \in \mathbb{R}^{64 \times 128}$,
$W_3 \in \mathbb{R}^{1 \times 64}$.  Trained jointly with BCE loss on
Polyvore outfit pairs.  Outfit-level scoring averages all
$\binom{n}{2}$ pairwise scores.

\subsection{Multi-Head Outfit Scoring}
\label{sec:mh_scoring}

The learned scorer replaces the rule-based four-level hierarchy with a
weighted combination of head-specific coherence signals:

\begin{align}
  S_{\text{total}} = \;&0.45 \cdot S_{\text{compat}} + 0.20 \cdot S_{\text{occasion}} \nonumber\\
                       &+ 0.20 \cdot S_{\text{style}} + 0.15 \cdot S_{\text{material}}
\label{eq:mh_score}
\end{align}

\noindent where:
\begin{itemize}[leftmargin=*,topsep=0pt,itemsep=1pt]
\item $S_{\text{compat}}$: mean pairwise compatibility score
  (Eq.~\ref{eq:compat})
\item $S_{\text{occasion}}$: mean cosine to centroid of items'
  \texttt{occasion} embeddings
\item $S_{\text{style}}$: mean cosine to centroid of items'
  \texttt{style} embeddings
\item $S_{\text{material}}$: harmony bonus ($+0.02$) if material
  embedding spread is in $[0.15, 0.45]$; penalty ($-0.02$) otherwise
\end{itemize}

\paragraph{Fallback.}
When multi-head embeddings are unavailable for any item (e.g., Shopify
catalog items not yet backfilled), the system falls back to the v1
rule-based scorer, ensuring graceful degradation.

\subsection{Object-Centric Wardrobe Ingestion}
\label{sec:ingestion}

Loom~v1 assumed one pre-cropped product photo per item.  Loom~v2 supports
three input modalities: single photo, video walkthrough, and live camera
scan.

\paragraph{Detection.}
YOLOv8-Nano~\citep{jocher2023ultralytics} detects clothing objects
(top, bottom, shoes, dress, bag) at confidence $\geq 0.4$.  On-device
inference runs at $\sim$15 FPS via ONNX Runtime Web; server-side
processes uploaded video at 3 FPS via ffmpeg extraction.

\paragraph{Tracking.}
ByteTrack~\citep{zhang2022bytetrack} assigns persistent object IDs across
frames using IoU-based bipartite matching with two-pass confidence
thresholds (high $\geq 0.5$, low $\geq 0.1$).  The Hungarian algorithm
(scipy, with greedy fallback) solves the assignment.  This converts a
stream of detections into a set of tracked objects, each with a temporal
crop buffer.

\paragraph{Best frame selection.}
For each tracked object, the system maintains a temporal buffer of the
top-$K$ ($K{=}5$) crops ranked by a composite quality score:
\begin{equation}
  Q = 0.5 \cdot \text{sharpness} + 0.3 \cdot \text{centeredness}
      + 0.2 \cdot \text{area}
\label{eq:quality}
\end{equation}

\noindent Sharpness is the variance of the Laplacian; centeredness
measures proximity to image center; area rewards larger bounding boxes.
Only the best crop per tracked object proceeds to segmentation and
embedding, reducing computation from $O(\text{frames})$ to
$O(\text{objects})$.

\paragraph{Temporal embedding aggregation.}
For video mode, each crop in the temporal buffer is segmented and embedded
independently, then combined via quality-weighted averaging:
\begin{equation}
  \mathbf{e}_{\text{agg}} = \text{normalize}\bigg(
    \sum_{k=1}^{K} Q_k \cdot \mathbf{e}_k
  \bigg)
\end{equation}

\paragraph{Async batch processing.}
Multi-item uploads are handled by an asynchronous ingestion worker with
concurrency 4, retry with exponential backoff, and a state machine
(pending $\to$ processing $\to$ ready/error) tracked per item.

\subsection{Segmentation-First Embedding}
\label{sec:segmentation}

All embeddings (FashionCLIP and DINOv2) are computed on
background-removed cutouts rather than raw photos.  We use
rembg~\citep{rembg} with the ISNet-general-use model:

\begin{enumerate}[leftmargin=*,topsep=2pt,itemsep=1pt]
\item Resize to max 1024px (speed)
\item Segment on natural background (better edge detection)
\item Composite RGBA result onto white background (JPEG, $Q{=}92$)
\end{enumerate}

The full original image is stored in Cloudinary for display; only the
cutout is used for embedding.  Segmentation failures gracefully fall back
to the original image.

\paragraph{Motivation.}
User-uploaded closet photos contain rooms, bodies, furniture, and
lighting artifacts.  Without segmentation, DINOv2 and FashionCLIP encode
\emph{scene} features rather than \emph{garment} features.  Segmentation
isolates the garment signal, improving embedding quality for both
retrieval and scoring.

\subsection{Cascade Retrieval with Sequential Conditioning}
\label{sec:cascade}

Loom~v1 retrieved candidates for all slots independently and scored
combinatorial candidates.  Loom~v2 introduces three-round cascade
retrieval that mirrors how stylists build outfits:

\paragraph{Round 1: Silhouette anchor.}
Bottom and shoes candidates are retrieved simultaneously.
\texttt{pick\_anchor\_pair()} evaluates $k{\times}k$ (bottom, shoes)
combinations and selects the pair with the tightest formality
coherence to the anchor item.

\paragraph{Round 2: Third piece.}
Layer and top candidates are retrieved with the chosen silhouette
pair as context---query construction includes the bottom and shoes
attributes, pushing retrieval toward pieces that complete the
emerging outfit.

\paragraph{Round 3: Finisher.}
Accessories are retrieved with the full outfit context
(anchor + silhouette + third piece), ensuring coordination.

\paragraph{Final scoring.}
All candidates from the three rounds are pooled per slot, deduplicated,
and passed to combinatorial generation (max 8 outfit candidates).  The
compatibility scorer selects the best outfit per direction.

\subsection{Multi-Head Retrieval}
\label{sec:mh_retrieval}

When multi-head embeddings are available, retrieval uses the
\texttt{compat\_embedding} column with cosine distance ($\Leftrightarrow$
operator) instead of the legacy FashionCLIP L2 distance ($\leftrightarrow$
operator).  The query vector is the anchor item's own
\texttt{compat\_embedding} (128-d), finding items that are
\emph{compatible with} the anchor rather than \emph{similar to} it.

Occasion filtering uses \texttt{occasion\_embedding} centroid coherence:
candidates are scored by cosine to the centroid of retrieved items'
occasion vectors, with a relative threshold (70\% of best score).
Direction reranking uses \texttt{style\_embedding} cosine to the anchor's
style vector.

Both mechanisms replace the hand-crafted vibe/anti-vibe prose strings and
per-occasion penalty coefficients from v1.

% ══════════════════════════════════════════════════════════════════════════════
\section{Experiments}
\label{sec:experiments}

\subsection{Experimental Setup}

We evaluate on the same catalog of 620 items (H\&M product data) used in
Loom~v1, with 50 randomly sampled anchor items (seed 42) generating
3 outfits each (150 outfits per configuration).  All items have been
backfilled with DINOv2 backbone and multi-head embeddings.

\subsection{Ablation Study}

\begin{table}[t]
\centering
\small
\caption{Multi-head ablation results (50 anchors $\times$ 3 directions,
single run).  \textbf{Bold} indicates best per column.}
\label{tab:mh_ablation}
\begin{tabular}{@{}lcccc@{}}
\toprule
Configuration & Score & Viol.\,\% & Slots & ms \\
\midrule
v1 full (rules)          &  0.011 & 8.0 & 4.68 & 6537 \\
\midrule
+ MH scoring only        & $-$0.018 & 10.7 & 4.68 & 6128 \\
+ MH occasion only       &  0.490 & \textbf{0.0} & 4.19 & 6037 \\
+ MH retrieval only      &  0.486 & \textbf{0.0} & 4.19 & 5564 \\
\midrule
MH full                  & \textbf{0.489} & \textbf{0.0} & 4.20 & \textbf{5090} \\
\bottomrule
\end{tabular}
\end{table}

\paragraph{Analysis.}

The multi-head full system achieves a 44$\times$ improvement in outfit
score (0.489 vs.\ 0.011) and completely eliminates hard violations (0\%
vs.\ 8.0\%).  Latency decreases 22\% (5090\,ms vs.\ 6537\,ms) because
the learned heads skip rule-based computation.

\paragraph{Component contributions.}
The largest gains come from \emph{retrieval} and \emph{occasion filtering}
rather than scoring.  Switching only occasion filtering to multi-head
raises the score from 0.011 to 0.490---the occasion head's centroid
coherence is dramatically more effective than vibe/anti-vibe prose
strings.  Similarly, replacing FashionCLIP L2 retrieval with
compat\_embedding cosine retrieval alone raises scores to 0.486.

\paragraph{Scoring-only is insufficient.}
Applying multi-head scoring alone ($-$0.018) slightly \emph{decreases}
performance.  This is expected: the compatibility scorer evaluates items
using their compat embeddings, but when retrieval still uses FashionCLIP
text queries, the retrieved candidates are optimized for text similarity
rather than compatibility---a distribution mismatch.  The scorer needs
compatible candidates to score, not merely similar ones.

\paragraph{Score interpretation caveat.}
The multi-head and rule-based systems use \emph{different scoring
functions}, so raw score magnitudes are not directly comparable.  What
\emph{is} comparable: violation rate (objective constraint checking),
slot diversity, and latency.  The 0\% violation rate under multi-head
scoring indicates that learned compatibility captures hard fashion
constraints without explicit rules.  A human preference study is needed
to validate subjective outfit quality (Section~\ref{sec:limitations}).

\subsection{Comparison with v1 Ablation}

\begin{table}[t]
\centering
\small
\caption{Loom v1 ablation (from~\citet{berlia2026loom}) alongside v2
multi-head results.  The $\Delta$ column shows change from the
respective system's full configuration.}
\label{tab:v1_comparison}
\begin{tabular}{@{}llcc@{}}
\toprule
System & Configuration & Score & Viol.\,\% \\
\midrule
v1 & Full (rules)           & 0.179 & 9.3 \\
v1 & $-$ Direction rerank   & 0.052 & 6.7 \\
v1 & Random retrieval       & 0.054 & 16.0 \\
\midrule
v2 & Full (multi-head)      & 0.489 & 0.0 \\
v2 & MH scoring only        & $-$0.018 & 10.7 \\
v2 & v1 full (rules, same run) & 0.011 & 8.0 \\
\bottomrule
\end{tabular}
\end{table}

Note that v1 and v2 ``full (rules)'' scores differ (0.179 vs.\ 0.011)
because the scoring function itself evolved between versions (flat
6-signal $\to$ 4-level hierarchy).  The relative ranking---multi-head
$\gg$ rules $>$ random---is consistent.

\subsection{Efficiency}

\begin{table}[t]
\centering
\small
\caption{End-to-end outfit generation latency (CPU, $N{=}50$ anchors).}
\label{tab:latency_v2}
\begin{tabular}{@{}lrr@{}}
\toprule
System & Mean (ms) & Speedup \\
\midrule
v1 rules (this run)    & 6,537 & --- \\
v2 multi-head full     & 5,090 & 1.28$\times$ \\
\bottomrule
\end{tabular}
\end{table}

The multi-head system is faster because:
(i)~compatibility scoring is a single forward pass through a small MLP
rather than hundreds of rule evaluations; and
(ii)~occasion filtering via centroid coherence is cheaper than embedding
each item against vibe/anti-vibe anchor vectors.

% ══════════════════════════════════════════════════════════════════════════════
\section{Ingestion Pipeline in Practice}
\label{sec:ingestion_practice}

The object-centric ingestion pipeline has been deployed for the personal
wardrobe application.  In practice:

\begin{itemize}[leftmargin=*,topsep=2pt,itemsep=1pt]
\item \textbf{Live camera scan}: The browser-side YOLO+ByteTrack system
  runs at $\sim$15 FPS on mobile devices via ONNX Runtime Web.  Stability
  thresholds (8 stable frames, minimum 10 track frames, quality $\geq 0.3$,
  1.5\,s cooldown) prevent redundant captures.
\item \textbf{Video upload}: Users record a closet walkthrough; the server
  extracts frames at 3 FPS, runs YOLO+ByteTrack, and processes the best
  crop per tracked object.  A 30-second video of 10 garments yields 10
  items processed in $\sim$45 seconds.
\item \textbf{Batch processing}: Multi-item uploads use an async worker
  with 4-way concurrency, retry with exponential backoff, and per-item
  state tracking (pending $\to$ processing $\to$ ready/error).
\end{itemize}

\paragraph{Reduction in computation.}
Tracking reduces embedding computation from $O(\text{frames} \times
\text{items})$ to $O(\text{items})$.  For a 30-second video at 3 FPS with
5 visible items, this reduces segmentation+embedding calls from
$90 \times 5 = 450$ to just 5.

% ══════════════════════════════════════════════════════════════════════════════
\section{Training Infrastructure}
\label{sec:training}

Head training operates on pre-computed backbone embeddings stored in HDF5,
enabling rapid iteration without re-running DINOv2:

\begin{enumerate}[leftmargin=*,topsep=2pt,itemsep=1pt]
\item \textbf{Backbone precomputation}: All training images are passed
  through DINOv2 once; 768-d CLS tokens are stored in
  \texttt{backbone\_embeddings.h5}.
\item \textbf{Pair construction}: Dataset-specific scripts
  (\texttt{prepare\_polyvore.py}, \texttt{prepare\_deepfashion.py},
  \texttt{prepare\_imaterialist.py}) generate positive/negative pair
  CSVs from outfit co-occurrence, attribute annotations, and the
  228-attribute iMaterialist taxonomy.
\item \textbf{Per-head training}: InfoNCE with $\tau{=}0.07$, batch 256,
  AdamW ($\text{lr}{=}10^{-3}$, $\text{wd}{=}10^{-2}$), 30 epochs.
  Compatibility head gets 40 epochs with hard negative mining (after
  epoch 10, 30\% of negatives are high-similarity non-compatible pairs).
\item \textbf{Scorer training}: Joint BCE on concatenated compat
  embeddings, same optimizer, 40 epochs.
\item \textbf{Evaluation}: Recall@$K$, compatibility AUC, and optional
  t-SNE/UMAP visualization via \texttt{evaluate\_heads.py}.
\end{enumerate}

Total trainable parameters: $\sim$1.2M (5 heads + scorer).  Training
completes in $<$30 minutes on a single GPU for all heads.

% ══════════════════════════════════════════════════════════════════════════════
\section{Database Schema Evolution}
\label{sec:schema}

The migration from v1 to v2 adds 6 vector columns per item table:

\begin{table}[t]
\centering
\small
\caption{Schema additions in migration 002.  HNSW indexes are created
on the three columns used in retrieval queries.}
\label{tab:schema}
\begin{tabular}{@{}lll@{}}
\toprule
Column & Dimensions & HNSW Index \\
\midrule
\texttt{backbone\_embedding} & 768 & No \\
\texttt{style\_embedding} & 128 & Yes \\
\texttt{fit\_embedding} & 128 & No \\
\texttt{material\_embedding} & 128 & No \\
\texttt{compat\_embedding} & 128 & Yes \\
\texttt{occasion\_embedding} & 128 & Yes \\
\bottomrule
\end{tabular}
\end{table}

The legacy \texttt{embedding vector(512)} column (FashionCLIP) is
retained for Shopify catalog items and fallback paths.  Total storage
overhead: 1,412 additional floats per item ($\sim$5.6\,KB at float32).

% ══════════════════════════════════════════════════════════════════════════════
\section{Discussion}
\label{sec:discussion}

\paragraph{Why DINOv2 over FashionCLIP?}
FashionCLIP's training objective (image--text alignment) produces
embeddings where ``looks like'' and ``goes with'' are conflated.  A red
sneaker and a red dress are close in FashionCLIP space because they share
visual attributes, but they are poor outfit partners.  DINOv2's
self-supervised objective produces features that capture \emph{visual
structure} without text-alignment bias.  Task-specific heads then learn
which aspects of visual structure matter for each downstream task.

\paragraph{Head independence.}
The five heads are trained independently, which is a deliberate design
choice.  Joint training (shared lower layers) would couple the
representations, making it harder to update individual heads (e.g.,
retraining the occasion head with new data) without affecting others.
The cost is slightly higher total parameters ($5 \times 200\text{K}$
vs.\ a shared trunk), which is negligible given the frozen backbone.

\paragraph{Scoring-retrieval alignment.}
The experiment confirms that scoring and retrieval must be aligned:
multi-head scoring with FashionCLIP retrieval ($-$0.018) performs
\emph{worse} than rule-based scoring with FashionCLIP retrieval (0.011).
This is because the compat scorer was trained on compat embeddings from
DINOv2, not on FashionCLIP similarity vectors---the input distribution
at inference doesn't match training.  When retrieval also switches to
compat embeddings, the scorer operates on in-distribution inputs and
performance jumps dramatically.

\paragraph{Removed: semantic material weight.}
The v1 paper's headline contribution---semantic material weight via CLIP
heavy/light context vectors---has been removed from the codebase.  The
original ablation showed it was overly conservative (removing it
\emph{increased} scores by $+$0.054).  The multi-head \texttt{material}
head provides a more nuanced material representation without hand-crafted
context strings.

% ══════════════════════════════════════════════════════════════════════════════
\section{Limitations}
\label{sec:limitations}

\begin{enumerate}[leftmargin=*,topsep=2pt,itemsep=1pt]
\item \textbf{No human evaluation.}  Outfit quality is ultimately
  subjective.  The 0\% violation rate and high coherence scores are
  encouraging but do not substitute for human preference judgments.
\item \textbf{Score non-comparability.}  Multi-head and rule-based scores
  use different scales and objectives; raw magnitudes cannot be compared
  across systems.
\item \textbf{Single catalog.}  All experiments use 620 H\&M items.
  Generalization to diverse brands, price ranges, and cultural contexts
  is untested.
\item \textbf{Shopify gap.}  Shopify catalog items use only FashionCLIP
  embeddings; multi-head benefits are limited to personal wardrobe and
  direct catalog modes.
\item \textbf{Head training data.}  Polyvore is no longer maintained;
  DeepFashion and iMaterialist have Western/Asian biases.  Broader
  training data would improve head generalization.
\end{enumerate}

% ══════════════════════════════════════════════════════════════════════════════
\section{Future Work}
\label{sec:future}

\begin{enumerate}[leftmargin=*,topsep=2pt,itemsep=1pt]
\item \textbf{Human preference study}: Side-by-side comparison of v1
  rule-based and v2 multi-head outfits rated by fashion-literate
  annotators.
\item \textbf{Shopify multi-head}: Extend DINOv2+head embeddings to
  Shopify catalog items via background backfill during catalog sync.
\item \textbf{Fine-tuning DINOv2}: Unfreeze the last 2 transformer blocks
  with low learning rate for fashion-specific feature adaptation.
\item \textbf{User preference heads}: Add a 6th ``personal'' head trained
  on individual user like/dislike feedback, replacing the EMA taste
  vector.
\item \textbf{Cross-catalog gap filling}: Use multi-head embeddings to
  identify wardrobe gaps and recommend retail items that would maximize
  outfit combinatorics.
\end{enumerate}

% ══════════════════════════════════════════════════════════════════════════════
\section{Conclusion}

We presented Loom~v2, which replaces a shared FashionCLIP embedding with
task-specific DINOv2 projection heads for outfit recommendation.  Five
independently trained heads (style, fit, material, compatibility,
occasion) and a learned pairwise scorer replace hundreds of lines of
rule-based scoring logic.  Object-centric ingestion via YOLOv8+ByteTrack
enables multi-item capture from video and live camera, transforming the
system from a static catalog tool into a persistent wardrobe understanding
platform.  Ablation experiments show the multi-head system achieves
44$\times$ higher outfit scores with zero constraint violations and 22\%
lower latency than its rule-based predecessor.

% ══════════════════════════════════════════════════════════════════════════════
\bibliographystyle{plainnat}
\begin{thebibliography}{25}

\bibitem[Berlia(2026a)]{berlia2026loom}
Berlia, A.
\newblock Loom: Hybrid retrieval-scoring outfit recommendation with semantic
material compatibility and occasion-aware embedding priors.
\newblock \emph{arXiv preprint}, 2026.

\bibitem[Berlia(2026b)]{berlia2026florence}
Berlia, A.
\newblock Fashion {Florence}: Fine-tuning {Florence-2} for structured fashion
attribute extraction.
\newblock \emph{arXiv preprint}, 2026.

\bibitem[Chen et~al.(2019)]{chen2019pog}
Chen, W., Huang, P., Xu, J., et~al.
\newblock {POG}: Personalized outfit generation for fashion recommendation at
{Alibaba iFashion}.
\newblock In \emph{KDD}, 2019.

\bibitem[Chia et~al.(2022)]{chia2022fashionclip}
Chia, P.~J., Attanasio, G., Bianchi, F., et~al.
\newblock {FashionCLIP}: Connecting language and images for product
representations.
\newblock \emph{arXiv preprint arXiv:2204.03972}, 2022.

\bibitem[Ge et~al.(2019)]{ge2019deepfashion2}
Ge, Y., Zhang, R., Wang, X., Tang, X., and Luo, P.
\newblock {DeepFashion2}: A versatile benchmark for detection, pose estimation,
segmentation and re-identification of clothing images.
\newblock In \emph{CVPR}, 2019.

\bibitem[Han et~al.(2019)]{han2019mvae}
Han, X., Song, J., and Wu, Z.
\newblock Learning fashion compatibility with context conditioning.
\newblock In \emph{AAAI}, 2019.

\bibitem[Jocher et~al.(2023)]{jocher2023ultralytics}
Jocher, G., Chaurasia, A., and Qiu, J.
\newblock Ultralytics {YOLOv8}.
\newblock \url{https://github.com/ultralytics/ultralytics}, 2023.

\bibitem[Lin et~al.(2020)]{lin2020csanet}
Lin, Y.-L., Tran, S., and Davis, L.~S.
\newblock Fashion outfit complementary item retrieval.
\newblock In \emph{CVPR}, 2020.

\bibitem[Oord et~al.(2018)]{oord2018infonce}
van~den Oord, A., Li, Y., and Vinyals, O.
\newblock Representation learning with contrastive predictive coding.
\newblock \emph{arXiv preprint arXiv:1807.03748}, 2018.

\bibitem[Oquab et~al.(2024)]{oquab2024dinov2}
Oquab, M., Darcet, T., Moutakanni, T., et~al.
\newblock {DINOv2}: Learning robust visual features without supervision.
\newblock \emph{TMLR}, 2024.

\bibitem[Radford et~al.(2021)]{radford2021clip}
Radford, A., Kim, J.~W., Hallacy, C., et~al.
\newblock Learning transferable visual models from natural language supervision.
\newblock In \emph{ICML}, 2021.

\bibitem[rembg()]{rembg}
rembg: Remove image background.
\newblock \url{https://github.com/danielgatis/rembg}.

\bibitem[Tan et~al.(2019)]{tan2019scenet}
Tan, R., Li, L., and Lin, D.
\newblock Learning conditional similarity networks for fashion recommendation.
\newblock In \emph{AAAI}, 2019.

\bibitem[Vasileva et~al.(2018)]{vasileva2018learning}
Vasileva, M., Plummer, B.~A., Dusber, K., et~al.
\newblock Learning type-aware embeddings for fashion compatibility.
\newblock In \emph{ECCV}, 2018.

\bibitem[Zhang et~al.(2022)]{zhang2022bytetrack}
Zhang, Y., Sun, P., Jiang, Y., et~al.
\newblock {ByteTrack}: Multi-object tracking by associating every detection
box.
\newblock In \emph{ECCV}, 2022.

\bibitem[Zheng et~al.(2018)]{zheng2018modanet}
Zheng, S., Yang, F., Kiapour, M.~H., and Piramuthu, R.
\newblock {ModaNet}: A large-scale street fashion dataset with polygon
annotations.
\newblock In \emph{ACM Multimedia}, 2018.

\end{thebibliography}

\end{document}
